% !TEX root = ../main.tex
%% Corrigé du sujet type bac n° 4 — valeurs vérifiées : verifications/sujets_bac.py

\corrigeSujet{4}

\begin{corrigeBac}{9}{L'eau oxygénée}
\partieExo{1}{Le titre d'une eau oxygénée}
\begin{questions}
\item Oxydation : $\ce{H2O2 -> O2 + 2H+ + 2e-}$ ; réduction : $\ce{H2O2 + 2H+ + 2e- -> 2H2O}$. En additionnant :
$\ce{2H2O2 -> 2H2O + O2}$. Le peroxyde d'hydrogène est à la fois l'\res{\text{oxydant}} et le \res{\text{réducteur}} : c'est une
dismutation. \pts{1 pt}
\item Un litre libère $\dfrac{10}{22{,}4}=\qty{0,446}{mol}$ de dioxygène ; il faut deux fois plus de peroxyde d'hydrogène :
$c=2\times0{,}446=\res{\qty{0,89}{mol.L^{-1}}}$. \pts{0,75 pt}
\end{questions}
\partieExo{2}{Suivi cinétique de la décomposition}
\begin{questions}[resume]
\item Initialement, $n_0=c_0\,V=0{,}20\times\num{20.0e-3}=\qty{4,0e-3}{mol}$. À l'avancement $x$ : $n(\ce{H2O2})=n_0-2x$ et $n(\ce{O2})=x$. Le dioxygène,
seul gaz formé, crée la surpression $\Delta P$ dans le volume $V_\mathrm{g}$ : $\Delta P\,V_\mathrm{g}=x\,R\,T$, d'où $\res{x=\dfrac{\Delta P\,V_\mathrm{g}}{R\,T}}$. \pts{1,25 pt}
\item $x_\mathrm{max}=\dfrac{n_0}{2}=\res{\qty{2,0e-3}{mol}}$, d'où $\Delta P_\mathrm{max}=\dfrac{\num{2.0e-3}\times8{,}314\times298}{\num{250e-6}}=\res{\qty{19,8}{kPa}}$, valeur vers
laquelle tend la courbe. \pts{0,75 pt}
\item Le temps de demi-réaction est la durée au bout de laquelle l'avancement atteint la moitié de sa valeur finale :
$\Delta P=\qty{9,9}{kPa}$ pour $\res{t_{1/2}=\qty{4,0}{min}}$. \pts{0,75 pt}
\item Les ions fer(III) sont dans la même phase que les réactifs : catalyse \res{\text{homogène}}. La vitesse diminue car la
concentration du réactif, facteur cinétique, diminue au cours du temps. \pts{0,75 pt}
\item Le programme calcule, à chaque date, $k=\dfrac1t\ln\dfrac{\Delta P_\mathrm{max}}{\Delta P_\mathrm{max}-\Delta P}$, valeur qui doit rester constante si
$\Delta P=\Delta P_\mathrm{max}\left(1-\mathrm{e}^{-kt}\right)$. Les valeurs affichées restent comprises entre 0,172 et \qty{0,175}{min^{-1}} : l'hypothèse
d'une loi d'\res{\text{ordre 1}} est validée, avec $\res{k\approx\qty{0,173}{min^{-1}}}$. \pts{1,5 pt}
\item $\dfrac{\ln2}{0{,}173}=\qty{4,0}{min}$ : on retrouve le temps de demi-réaction. Vitesse initiale :
$v_0=k\,c_0=0{,}173\times0{,}20=\res{\qty{0,035}{mol.L^{-1}.min^{-1}}}$. \pts{0,75 pt}
\item La catalase est un catalyseur \res{\text{enzymatique}} très efficace, et la température du corps, \qty{37}{\degreeCelsius}, est plus
élevée que celle du laboratoire. \pts{0,75 pt}
\item La lumière et la chaleur accélèrent la décomposition : le flacon opaque et la conservation au frais la
\res{\text{ralentissent}}, et l'eau oxygénée garde son titre plus longtemps. \pts{0,75 pt}
\end{questions}
\schema[corrige]{bac-s4-pression}
\end{corrigeBac}

\begin{corrigeBac}{6}{Une lunette astronomique}
\begin{questions}
\item Une lunette afocale donne d'un objet à l'infini une image à l'infini, observable sans accommoder : le foyer objet
$F_2$ de l'oculaire est confondu avec le foyer image $F'_1$ de l'objectif. Longueur : $f'_1+f'_2=900+10=\res{\qty{910}{mm}}$. \pts{1 pt}
\item Le rayon qui passe par $O_1$ n'est pas dévié ; tous les rayons du faisceau convergent en $B_1$, dans le plan focal
image de l'objectif. $B_1$ étant dans le plan focal objet de l'oculaire, le faisceau émerge parallèle à la droite $(B_1O_2)$
(voir le schéma). \pts{1,25 pt}
\item $G=\dfrac{\alpha'}{\alpha}$, rapport des angles sous lesquels on voit l'astre avec et sans la lunette. Avec $A_1B_1$ l'image
intermédiaire : $\alpha\approx\dfrac{A_1B_1}{f'_1}$ et $\alpha'\approx\dfrac{A_1B_1}{f'_2}$, d'où $\res{G=\dfrac{f'_1}{f'_2}}$. \pts{1 pt}
\item $G=\dfrac{900}{10}=\res{90}$ et $G=\dfrac{900}{25}=\res{36}$. \pts{0,5 pt}
\item $\alpha'=90\times45''=4050''$, soit $\res{\ang{1,1}}$ : Jupiter paraît deux fois plus grande que la Lune à l'œil nu. \pts{0,75 pt}
\item $\theta=\dfrac{\num{550e-9}}{0{,}070}=\qty{7,9e-6}{rad}$, soit $\res{1{,}6''}$, inférieur à $4''$ : la lunette sépare les bandes. Avec $G=90$, une bande
est vue sous $90\times4''=360''=6'$, bien plus que $1'$ : l'œil la distingue. L'oculaire de \qty{4}{mm} ($G=225$) agrandirait
l'image sans révéler de nouveaux détails : le plus petit détail séparé par l'objectif, $1{,}6''$, est déjà vu sous
$90\times1{,}6''=144''$, soit $2{,}4'$, avec $G=90$. L'image serait seulement plus grande et plus sombre. \pts{1,5 pt}
\end{questions}
\schema[corrige]{bac-s4-lunette}
\end{corrigeBac}

\begin{corrigeBac}{5}{Le château d'eau}
\begin{questions}
\item Loi de la statique des fluides : $P=P_\mathrm{atm}+\rho\,g\,h=\num{1.013e5}+1000\times9{,}81\times35=\res{\qty{4,4e5}{Pa}}$, soit une
surpression de \qty{3,4}{bar}. \pts{1 pt}
\item La pression au robinet augmente avec la différence de hauteur : la hauteur du château d'eau fournit la pression du
réseau. Au dernier étage, à \qty{40}{m}, plus haut que la surface libre, l'eau \res{\text{n'arrive pas}} : il faut un surpresseur. \pts{0,75 pt}
\item Entre la surface libre ($P_\mathrm{atm}$, vitesse négligeable, $z=h$) et la sortie du robinet ($P_\mathrm{atm}$, vitesse $v$, $z=0$) :
$\rho\,g\,h=\frac12\,\rho\,v^2$, d'où $v=\sqrt{2\,g\,h}=\sqrt{2\times9{,}81\times35}=\res{\qty{26}{m.s^{-1}}}$. Hypothèses : fluide parfait (sans
frottement), incompressible, écoulement permanent. \pts{1 pt}
\item $D_V=S\,v=\pi\times(\num{6.0e-3})^2\times26{,}2=\qty{2,96e-3}{m^3.s^{-1}}$, soit $\res{\qty{178}{L.min^{-1}}}$. \pts{0,75 pt}
\item $v=\dfrac{D_V}{S}=\dfrac{\num{12e-3}/60}{\num{1.13e-4}}=\res{\qty{1,8}{m.s^{-1}}}$, quinze fois moins que prévu : le modèle du fluide parfait ne
convient pas, les frottements dans les canalisations (pertes de charge) dissipent l'essentiel de l'énergie. Baignoire :
$\dfrac{150}{12}=\res{\qty{12,5}{min}}$. \pts{1,5 pt}
\end{questions}
\end{corrigeBac}
